% YÖK-uyumlu Türkçe tez iskeleti — scholar-kit
% Derleme: bash scripts/latex_build.sh templates/thesis-tr/main.tex build
% Kurum kılavuzu templates/institution/ altında kayıtlıysa onun kuralları önceliklidir
% (TEMPLATE SUPREMACY — bkz. thesis-structure skill, reference/yok-tr.md).
\documentclass[12pt,a4paper,oneside]{report}

% --- Dil ve yazı tipi ---
\usepackage[utf8]{inputenc}      % pdflatex için; xelatex kullanılıyorsa fontspec'e geçin
\usepackage[T1]{fontenc}
\usepackage[turkish,shorthands=:!]{babel}  % shorthands çakışmalarını önler
\usepackage{times}               % kurum kılavuzu farklı font isterse değiştirin

% --- Sayfa düzeni (tipik YÖK: sol 3.5cm, diğerleri 2.5cm) ---
\usepackage[left=3.5cm,right=2.5cm,top=2.5cm,bottom=2.5cm]{geometry}
\usepackage{setspace}\onehalfspacing

% --- Kaynakça: tek kaynak bib/references.bib; stil PROJECT.md'den ---
\usepackage[backend=biber,style=apa,language=turkish]{biblatex}
\addbibresource{../../bib/references.bib}

% --- Şekil/tablo/çapraz referans ---
\usepackage{graphicx,booktabs,caption,subcaption}
\usepackage[hidelinks]{hyperref}
\usepackage[capitalise]{cleveref}  % cleveref has no "turkish" language option — see below
% cleveref ships no Turkish translations; define them manually (bkz. thesis-structure skill)
\crefname{chapter}{bölüm}{bölüm}     \Crefname{chapter}{Bölüm}{Bölüm}
\crefname{section}{kısım}{kısım}     \Crefname{section}{Kısım}{Kısım}
\crefname{figure}{şekil}{şekil}      \Crefname{figure}{Şekil}{Şekil}
\crefname{table}{tablo}{tablo}       \Crefname{table}{Tablo}{Tablo}
\crefname{equation}{denklem}{denklem} \Crefname{equation}{Denklem}{Denklem}
\crefname{page}{sayfa}{sayfa}        \Crefname{page}{Sayfa}{Sayfa}

\begin{document}

% ===== ÖN SAYFALAR (YÖK sırası — yok-tr.md) =====
% 1. İç kapak
% İç kapak — kurum kılavuzundaki düzeni birebir uygulayın (templates/institution/)
\begin{titlepage}\centering
{\large ÜNİVERSİTE ADI}\\ {\large ENSTİTÜ ADI}\\[3cm]
{\LARGE\bfseries TEZ BAŞLIĞI}\\[2cm]
{\large Yazar Adı SOYADI}\\[1cm]
{\large YÜKSEK LİSANS TEZİ / DOKTORA TEZİ}\\ {\large Anabilim Dalı}\\[2cm]
{\large Danışman: Unvan Ad SOYAD}\\[1cm]
{\large Şehir, Yıl}
\end{titlepage}

% 2. Onay sayfası
% Onay sayfası — enstitü şablonundan alınır; jüri imza alanları
\chapter*{Onay Sayfası}\addcontentsline{toc}{chapter}{Onay Sayfası}
[Enstitü onay sayfası şablonu buraya]

% 3. Etik beyan
\chapter*{Etik Beyan}\addcontentsline{toc}{chapter}{Etik Beyan}
Bu tezdeki bütün bilgilerin akademik kurallara ve etik ilkelere uygun olarak
elde edildiğini ve sunulduğunu; bana ait olmayan her türlü ifade ve bilginin
kaynağına eksiksiz atıf yapıldığını beyan ederim.

% 4. Teşekkür
\chapter*{Teşekkür}\addcontentsline{toc}{chapter}{Teşekkür}
[Teşekkür metni]

% 5. Özet (TR) — Abstract ile yapısal eşdeğer olmalı
% Özet — Abstract ile YAPISAL EŞDEĞER olmalı (aynı paragraf yapısı, aynı iddialar,
% aynı anahtar kelimeler) — /submit-check bunu style-auditor ile denetler.
\chapter*{Özet}\addcontentsline{toc}{chapter}{Özet}
[Özet metni — drafts/ozet.md kaynağından]
\par\noindent\textbf{Anahtar Kelimeler:} kelime1, kelime2, kelime3, kelime4, kelime5

% 6. Abstract (EN)
\chapter*{Abstract}\addcontentsline{toc}{chapter}{Abstract}
[Abstract text — structurally equivalent to Özet]
\par\noindent\textbf{Keywords:} keyword1, keyword2, keyword3, keyword4, keyword5

% 7-9. İçindekiler, tablo/şekil/simge listeleri
\tableofcontents
\listoftables
\listoffigures
% \input{frontmatter/simgeler}

% ===== BÖLÜMLER (5 bölümlü klasik yapı — thesis-structure skill) =====
\chapter{Giriş}
% drafts/01-giris.md içeriği /draft ile üretilir; format-engineer buraya taşır
          % RQ'lar burada tanımlanır (RQ1, RQ2, ...)
\chapter{Literatür Taraması}
% drafts/02-literatur.md içeriği /draft ile üretilir; format-engineer buraya taşır

\chapter{Yöntem}
% drafts/03-yontem.md içeriği /draft ile üretilir; format-engineer buraya taşır

\chapter{Bulgular}
% drafts/04-bulgular.md içeriği /draft ile üretilir; format-engineer buraya taşır
       % Başlıklar RQ'larla 1:1 eşleşmeli (--rq-trace)
\chapter{Tartışma ve Sonuç}
% drafts/05-sonuc.md içeriği /draft ile üretilir; format-engineer buraya taşır


\printbibliography[title={Kaynakça}]

% \appendix
% \input{appendices/ek-a}

\end{document}
